\section{Conclusion and Future Work}
In this paper, we proposed a robotic decision-making framework in which human feedback is treated as targeted state verification rather than action delegation. The proposed method maintains a belief over symbolic possible worlds within a POMDP-based planning framework and quantifies uncertainty through a probability distribution over these worlds. Based on this belief, the robot decides when human input is needed and asks the user to verify a planning-relevant predicate fact. The answer filters inconsistent possible worlds, allowing the robot to preserve responsibility for action selection while using human input to resolve state ambiguity.

In particular, we use entropy-based confidence to determine when to issue queries and select the hypothesis expected to reduce belief ambiguity. Experimental results show that selective state verification can maintain task performance while reducing interaction relative to dense supervision. The comparison with action-level querying highlights a complementary trade-off: action guidance can reduce query frequency and shorten execution in some settings, whereas state verification directly resolves the symbolic uncertainty used by the planner. The preliminary user study further suggests that state-level queries can reduce interaction effort and workload, while also revealing domain-dependent effects on situation awareness.

These results point to several extensions. First, richer feedback channels can build on the current binary hypothesis interface by allowing users to provide corrected symbolic facts or natural-language clarifications. Second, scalable frontier construction can extend the approach to larger domains through pruning, abstraction, or learned proposal distributions. Third, user-aware query selection can combine robot belief uncertainty with estimates of what users can judge from the interface. Finally, the framework can be extended to multi-robot settings, where multiple agents share uncertainty and coordinate when and whom to query. These extensions preserve the central idea of the framework: human feedback is most useful when it is requested at belief-relevant moments and expressed as information that helps the robot resolve state uncertainty without outsourcing its planning decisions.
